
% SGA 3，第 VII 篇_B；英文重建本第一轮的中文译本。
% 范围： 节 2.6.2 全文 依据文献局部页 61--62,
% 止于下节之前： 节 2.6.3.
% 合并阅读本页： 574--575. 重版印刷页： 563--564.
% 范围内原版边码： 528.
% 范围内编者注： 110.
% 依据： Polo--Gille Exp7B-13oct24.pdf,
% SHA-256 D303690E0195654BBE7C3DAA5D89EA23E24CD709A3575217234AC19AE8836694.
% 参照了 Jacob Reinhold 的 CC BY 4.0 Markdown，作为对照来源。
% 本组件无图；所有数学内容均为原生 LaTeX。
\subsubsection*{2.6.2}
\phantomsection
\label{sga3:VIIB:2.6.2}
反过来，每个 \(\mathbf O_k\) 上的 Lie 代数 \(\mathbf L\) 都定义一个 \(k\)-群函子。事实上，令 \(U(\mathbf L)\) 表示将 \(\operatorname{Alf}/k\) 的每个对象 \(C\) 映为 \(C\)-Lie 代数 \(\mathbf L(C)\) 的包络代数
\[
  U\bigl(\mathbf L(C)\bigr)
\]
的函子。由 VII\textsubscript{A}，\SGARefLink{sga3:VIIA:3.2.2}{3.2.2}，\(U(\mathbf L)\) 是 \(\mathbf O_k\) 上的双代数，故由 \SGARefLink{sga3:VIIB:2.2}{2.2} 确定一个 \(k\)-群函子
\[
  \operatorname{Spf}^{*}U(\mathbf L),
\]
此后将它记为 \(\mathscr G(\mathbf L)\)。于是 \(\mathscr G(\mathbf L)(C)\) 是由增广为 \(1\) 且满足
\[
  \Delta_{U(\mathbf L(C))}(z)=z\otimes z.
\]
的元素 \(z\in U(\mathbf L(C))\) 构成的群。
当 \(\mathbf L\) 在 \(\mathbf O_k\) 上平坦时，有如下命题。
\medskip
\noindent\textbf{命题。}\footnote{%
\label{sga3:VIIB:nde:110}编者将\textup{\SGARefLink{sga3:VIIB:2.6.2:prop:item:i}{(i)}} 和\textup{\SGARefLink{sga3:VIIB:2.6.2:prop:item:ii}{(ii)}} 分项列出，并添加了第\textup{(iii)} 项，后者将用于 \SGARefLink{sga3:VIIB:2.6.3:prop}{2.6.3} 和 \SGARefLink{sga3:VIIB:3.3.2}{3.3.2}。}
\phantomsection
\label{sga3:VIIB:2.6.2:prop}
{\itshape 设 \(\mathbf L\) 是 \(\mathbf O_k\) 上的平坦 Lie 代数。
\begin{enumerate}
\item[\textup{\SGAInlineRefTarget{sga3:VIIB:2.6.2:prop:item:i}(i)}]
\(\mathscr G(\mathbf L)\) 是在 \(k\) 上拓扑平坦的形式群，其 \(\mathbf O_k\)-双代数为 \(U(\mathbf L)\)。
\item[\textup{\SGAInlineRefTarget{sga3:VIIB:2.6.2:prop:item:ii}(ii)}] \(\mathscr G(\mathbf L)\) 是无穷小的。
\item[\textup{\SGAInlineRefTarget{sga3:VIIB:2.6.2:prop:item:iii}(iii)}]
对 \(\operatorname{Alf}/k\) 的每个对象 \(C\)，\(\operatorname{Lie}(\mathscr G(\mathbf L))(C)\) 可与 \(U(\mathbf L(C))\) 的本原元素集合
\[
  \operatorname{Prim}U(\mathbf L(C))
  =
  \left\{
    x\in U(\mathbf L(C))
    \,\middle|\,
    \varepsilon(x)=0,\quad
    \Delta(x)=x\otimes1+1\otimes x
  \right\}
\]
等同。特别地，有 \(\mathbf O_k\) 上的 Lie 代数的自然态射
\[
  \tau_{\mathbf L}:
  \mathbf L\longrightarrow
  \operatorname{Lie}\bigl(\mathscr G(\mathbf L)\bigr).
\]
\end{enumerate}}
\textit{证明。}
\(\mathbf L\) 在 \(\mathbf O_k\) 上平坦，意即对 \(\operatorname{Alf}/k\) 中的每个态射 \(C\to D\)，
\[
  \mathbf L(D)=\mathbf L(C)\otimes_C D,
\]
且对 \(C\) 的每个局部分量 \(C'\)，\(C'\)-模 \(\mathbf L(C')\) 是自由的。第一个条件结合张量积和包络代数的泛性质，给出
\[
  U\bigl(\mathbf L(D)\bigr)
  =
  U\bigl(\mathbf L(C)\bigr)\otimes_C D.
\]
第二个条件和 Poincaré--Birkhoff--Witt 定理\textup{（参见 Bourbaki，《Lie 群与 Lie 代数》〔\emph{Groupes et algèbres de Lie}〕，I，2.7）}表明 \(U(\mathbf L(C'))\) 是自由 \(C'\)-模。因此，双代数 \(U(\mathbf L)\) 在 \(\mathbf O_k\) 上平坦，证明了\textup{\SGARefLink{sga3:VIIB:2.6.2:prop:item:i}{(i)}}。
为证明\textup{\SGARefLink{sga3:VIIB:2.6.2:prop:item:ii}{(ii)}} 和\textup{\SGARefLink{sga3:VIIB:2.6.2:prop:item:iii}{(iii)}}，可设 \(k\) 是 Artin 环。置
\[
  L=\mathbf L(k),\qquad
  U=U(L),\qquad
  U_0=k\cdot1_U,
\]
并令 \(U^+\) 为 \(U\) 中由 \(L\) 的像生成的双边理想。对每个 \(n\geq1\)，置
\[
  U_n
  =
  \left\{
    u\in U
    \,\middle|\,
    \Delta_U(u)-u\otimes1\in U_{n-1}\otimes U^+
  \right\}.
\]
由 \SGARefLink{sga3:VIIB:1.3.6}{1.3.6}，只须证明 \(U\) 是各 \(U_n\) 的并。将 \(L\) 与其在 \(U\) 中的典范像等同，则 \(L\subset U_1\)。若 \(x_1,\ldots,x_n\in L\)，则
\[
  \Delta_U(x_1\cdots x_n)
  =
  (x_1\otimes1+1\otimes x_1)\cdots
  (x_n\otimes1+1\otimes x_n).
\]
对 \(n\) 归纳，即得 \(x_1\cdots x_n\in U_n\)，从而
\[
  U=\bigcup_n U_n.
\]
这就证明了\textup{\SGARefLink{sga3:VIIB:2.6.2:prop:item:ii}{(ii)}}。
最后，令
\[
  D=k[d]/(d^2)
\]
为 \(k\) 上的对偶数代数。由假设，
\[
  \mathbf L(D)\simeq L\otimes_kD,
  \qquad
  U(\mathbf L(D))\simeq U\otimes_kD.
\]
由此可得 \(\operatorname{Lie}(\mathscr G(\mathbf L))(k)\) 可与满足 \(\varepsilon(z)=1\) 和 \(\Delta(z)=z\otimes z\) 的元素
\[
  z=1+xd\in U\oplus Ud
  \qquad (x\in U)
\]
的集合等同。这些条件等价于
\[
  \varepsilon(x)=0,
  \qquad
  \Delta(x)=x\otimes1+1\otimes x,
\]
即 \(x\in\operatorname{Prim}U\)。特别地，
\[
  \tau_{\mathbf L}:x\longmapsto1+dx
\]
是从 \(\mathbf L\) 到 \(\operatorname{Lie}(\mathscr G(\mathbf L))\) 的 \(\mathbf O_k\) 上 Lie 代数态射。
